\hypertarget{the-disassembler-dis}{%
\chapter{\texorpdfstring{The Disassembler
(\texttt{dis})}{The Disassembler (dis)}}\label{the-disassembler-dis}}

To reach ``Godhood,'' you must be able to read the machine code of the
Python Virtual Machine: \textbf{Bytecode}.

\hypertarget{the-python-vm-a-stack-machine}{%
\subsection{62.1 The Python VM: A Stack
Machine}\label{the-python-vm-a-stack-machine}}

The CPython VM is a \textbf{Stack Machine}. Operations push values onto
a stack and pop them to perform calculations.

\hypertarget{dissecting-an-operation}{%
\subsubsection{1. Dissecting an
Operation}\label{dissecting-an-operation}}

\begin{lstlisting}[language=Python]
def add(a, b):
    return a + b

import dis
dis.dis(add)
\end{lstlisting}

\emph{Output:}

\begin{lstlisting}
  2           0 LOAD_FAST                0 (a)
              2 LOAD_FAST                1 (b)
              4 BINARY_ADD
              6 RETURN_VALUE
\end{lstlisting}

\begin{itemize}
\tightlist
\item
  \textbf{\passthrough{\lstinline!LOAD\_FAST!}}: Pushes the value of a
  local variable onto the stack.
\item
  \textbf{\passthrough{\lstinline!BINARY\_ADD!}}: Pops the top two
  values, adds them (using the
  \passthrough{\lstinline!tp\_as\_number->nb\_add!} C slot), and pushes
  the result back.
\item
  \textbf{\passthrough{\lstinline!RETURN\_VALUE!}}: Pops the top value
  and returns it to the caller.
\end{itemize}

\hypertarget{bytecode-specialization-python-3.11}{%
\subsection{62.2 Bytecode Specialization (Python
3.11+)}\label{bytecode-specialization-python-3.11}}

In modern Python, you may see \passthrough{\lstinline!RESUME!} or
``Specialized'' opcodes like
\passthrough{\lstinline!BINARY\_OP\_ADD\_INT!}. * \textbf{Inline
Caching}: If the VM sees that a specific
\passthrough{\lstinline!BINARY\_ADD!} is always adding two integers, it
replaces the generic opcode with a specialized version that skips the
type-checking overhead, resulting in significant speedups (as discussed
in Chapter 17).

\begin{center}\rule{0.5\linewidth}{0.5pt}\end{center}
